\subsection{The Klein bottle quasivariety: Problem 17.33}
\label{cand:17.33}
\problemstatement{17.33}

Let \(K=\Z\rtimes\Z\) with
\((u,k)(v,l)=(u+(-1)^kv,k+l)\), the Klein bottle group.
Its quasivariety \(q(K)\) cannot be axiomatized by
quasiidentities with any fixed finite bound on the number
of variables.

For \(n\geq2\), put \(E=\F_2^n\), and index \(m=n(2^n-1)\)
integer coordinates by pairs \((\chi,j)\), where
\(0\ne\chi\in E^*\) and \(1\leq j\leq n\).
Let \(D_e\) multiply coordinate \((\chi,j)\) by
\((-1)^{\chi(e)}\). Define
\[
 \Gamma_n=\{(e,c)\in E\times\Z^m:
              c_{\chi,j}\equiv e_j\pmod2\},\qquad
 (e,c)(f,d)=(e+f,c+D_ed).
\]
Parity makes this a subgroup of the affine semidirect product.
Its translation subgroup \(N=\{(0,2z):z\in\Z^m\}\)
has quotient \(E\).

The group is torsion-free. If \(e\ne0\), choose a nonzero
\(\chi\) vanishing on \(e\), possible because \(n\geq2\),
and \(j\) with \(e_j=1\). Then the \((\chi,j)\)-coordinate
of \((e,c)^2=(0,(1+D_e)c)\) is nonzero. A nonidentity
translation also has nonzero square, and any torsion element
would have square one since its square lies in the
torsion-free group \(N\). On the other hand, the
abelianization of \(\Gamma_n\) is finite. Every translation
basis vector is inverted by some lift of \(e_i\), since
\(\chi_i=1\) for some \(i\), so \(2N\leq\Gamma_n'\).
The quotient by \(2N\) is finite.

Every homomorphism \(\Gamma_n\to K\) is trivial: composition
with \(K\to\Z\) is trivial by finite abelianization, and
the image is then in the translation subgroup \(\Z\), where
the same argument applies again.

In contrast, the inverse image \(\Gamma_L\) of every proper
subspace \(L<E\) embeds in \(K^m\). Choose \(0\ne\psi\in E^*\)
vanishing on \(L\), and define, with coefficients
\(\chi_j\in\{0,1\}\),
\[
 \ell_\chi(e,c)=\sum_j\chi_j c_{\psi,j},\qquad
 F_{\chi,j}(e,c)=(c_{\chi,j},\ell_\chi(e,c)).
\]
The integer-valued \(\ell_\chi\) is additive on \(\Gamma_L\),
because \(D_e\) fixes the \(\psi\)-coordinates for \(e\in L\),
and it has parity \(\chi(e)\). Thus every \(F_{\chi,j}\)
is a homomorphism into \(K\). Their product is injective:
all its first coordinates determine \(c\), and the parities
of \(c\) determine \(e\). In particular, every subgroup
generated by fewer than \(n\) elements belongs to \(q(K)\).

It remains to justify that \(\Gamma_n\notin q(K)\), without
assuming a residual characterization that might require
additional hypotheses. The group is finitely presented.
Explicit generators are translations \(t_{\chi,j}\) and
lifts \(g_i=(e_i,c_i)\) with
\((c_i)_{\chi,j}=\delta_{ij}\). A finite presentation is
given by commuting translations and the relations
\[
 \begin{split}
 g_it_{\chi,j}g_i^{-1}&=t_{\chi,j}^{(-1)^{\chi_i}},\\
 g_i^2&=\prod_{\chi:\chi_i=0}t_{\chi,i},\\
 g_ig_j&=\left(\prod_\chi
          t_{\chi,i}^{\chi_j}t_{\chi,j}^{-\chi_i}\right)g_jg_i
          \quad(i<j).
 \end{split}
\]
These hold by the affine multiplication law. They collect
every word to a translation monomial followed by
\(g_1^{\epsilon_1}\cdots g_n^{\epsilon_n}\),
\(\epsilon_i\in\{0,1\}\). Its affine image is the identity
only if all \(\epsilon_i\) and then all translation
exponents vanish, proving the presentation.

For such a presentation, the quasiidentity asserting that
its finitely many relations imply that a chosen nontrivial
generator is one holds in \(K\), since every homomorphism
\(\Gamma_n\to K\) is trivial. It fails in \(\Gamma_n\).
Thus \(\Gamma_n\notin q(K)\). If quasiidentities in at most
\(r\) variables axiomatized \(q(K)\), take \(n>r\).
Every tuple of at most \(r\) elements of \(\Gamma_n\)
lies in a subgroup of \(q(K)\), so all those
quasiidentities hold in \(\Gamma_n\), a contradiction.
Source: \artifact{research/17.33-proof.md}.

